\section{问题重述与任务分析}
\subsection{问题背景}
机器狗携带测向传感器和近距离光学清除装置，从目标区域中心出发，按照接口反馈逐步寻找无线电干扰源。干扰源的位置、频道和总数都不预先公开；一次动作只能针对一个频道和一个位置，接口返回当前频道的 \texttt{direction}、\texttt{near} 或 \texttt{no\_signal}。测向角存在不超过 $1^\circ$ 的误差，机器狗的运动速度为 $5\,\mathrm{m/s}$，测量、换频、移动和清除都计入同一虚拟时钟。源进入 $20\,\mathrm m$ 范围后才能尝试清除，因而“发现了方向”与“已经可以安全清除”是两个不同的停止条件。题目背景和接口约束以竞赛官网公布的规则为准\cite{cumcm}。

本题的困难来自信息的不对称。程序知道自己走过的坐标和接口反馈，却不能读取隐藏源坐标，也不能把一次没有信号解释为频道为空。第三问的源均全向，接收半径决定可见性；第四问的定向源还只在一个未知闭半平面内发射，同一位置换一个测量站可能从有信号变为无信号。策略必须在不知道真值的情况下逐步缩小可能区域，并且在结束前给出覆盖和清除两类可核查证据。

第三问的源均为全向源，重点是用有限站点完成全域发现，再根据多次测向收缩位置区域。第四问增加未知发射方向：检测点即使距离足够近，也可能位于发射半平面之外，因此“无信号”不能直接当作该频道不存在。该变化使覆盖证书和定位策略必须同时考虑。

\subsection{任务的输入、输出与评价}
\begin{table}[H]\centering\small
\caption{题目要求、输出和本文方法}\label{tab:q1-tasks}
\begin{tabular}{p{0.12\linewidth}p{0.35\linewidth}p{0.39\linewidth}}
\toprule
问题 & 输入与约束 & 输出与评价 \\
\midrule
1 & 多个检测位置、示向度和误差界 & 可行定位区域、直径、最小包围圆及一次清除判据 \\
2 & 首次检测位置、示向度及候选移动位置 & 可接收候选集、候选点排序和后验定位误差 \\
3 & 串行反馈、全向源、区域覆盖要求 & 全部源的发现、定位、清除策略及虚拟时间 \\
4 & 问题3信息，加未知源类型和发射方向 & 面向定向盲区的覆盖策略、清除证书和鲁棒性评价 \\
\bottomrule
\end{tabular}
\end{table}

表~\ref{tab:q1-tasks} 中的“输入”均指题目给出的规则和当前接口反馈，而非合成环境中的隐藏真值。程序把每个动作的请求、响应和虚拟时钟写入逐动作日志；论文中的数值只从这些日志或固定种子合成运行中读取。评价先看是否完成覆盖并清除全部源，再比较虚拟时间；若两者冲突，完整性优先。反馈驱动的决策循环见图~\ref{fig:q1-workflow}。

\subsection{建模主线}
我们将每个频道的状态写成 $(\mathcal P_k,\,z_k,\,c_k)$：$\mathcal P_k$ 是由已有反馈保留下来的可能位置集合，$z_k$ 表示是否已发现，$c_k$ 表示是否已清除。一次 \texttt{direction} 反馈把一个带角度误差的楔形约束加入 $\mathcal P_k$；\texttt{near} 直接触发同位置清除；\texttt{no\_signal} 只更新测量记录，不擅自删除候选位置。每一次动作先通过协议客户端取得反馈，再更新对应状态，策略不读取隐藏源坐标。

决策目标按词典序排列：先满足覆盖和清除证书，再比较总虚拟时间、移动距离、测量次数和兜底次数。问题四的覆盖阶段和末端阶段因此承担不同任务：前者完成对未知朝向的发现证书并积累方向信息，后者在可能区域已经收缩后集中清除。这样的排序避免把“较快但漏源”的结果误当成优解，也为后面的算法比较提供统一口径。

\begin{figure}[H]\centering
\begin{tikzpicture}[>=Latex,node distance=10mm]
\node[draw,rounded corners,fill=gray!10,minimum width=30mm,minimum height=9mm] (a) {题面/接口反馈};
\node[draw,rounded corners,fill=gray!10,minimum width=30mm,minimum height=9mm,right=of a] (b) {集合状态更新};
\node[draw,rounded corners,fill=gray!10,minimum width=30mm,minimum height=9mm,right=of b] (c) {下一停靠点};
\node[draw,rounded corners,fill=gray!10,minimum width=30mm,minimum height=9mm,right=of c] (d) {清除证书};
\draw[->] (a)--(b);\draw[->] (b)--(c);\draw[->] (c)--(d);\draw[->,rounded corners] (d) |- ++(0,-12mm) -| (a);
\end{tikzpicture}
\caption{反馈驱动的在线决策循环}\label{fig:q1-workflow}
\source{本文算法流程；不含隐藏真值。}
\end{figure}
